% Options for packages loaded elsewhere
\PassOptionsToPackage{unicode}{hyperref}
\PassOptionsToPackage{hyphens}{url}
%
\documentclass[
]{article}
\usepackage{amsmath,amssymb}
\usepackage{iftex}
\ifPDFTeX
  \usepackage[T1]{fontenc}
  \usepackage[utf8]{inputenc}
  \usepackage{textcomp} % provide euro and other symbols
\else % if luatex or xetex
  \usepackage{unicode-math} % this also loads fontspec
  \defaultfontfeatures{Scale=MatchLowercase}
  \defaultfontfeatures[\rmfamily]{Ligatures=TeX,Scale=1}
\fi
\usepackage{lmodern}
\ifPDFTeX\else
  % xetex/luatex font selection
\fi
% Use upquote if available, for straight quotes in verbatim environments
\IfFileExists{upquote.sty}{\usepackage{upquote}}{}
\IfFileExists{microtype.sty}{% use microtype if available
  \usepackage[]{microtype}
  \UseMicrotypeSet[protrusion]{basicmath} % disable protrusion for tt fonts
}{}
\makeatletter
\@ifundefined{KOMAClassName}{% if non-KOMA class
  \IfFileExists{parskip.sty}{%
    \usepackage{parskip}
  }{% else
    \setlength{\parindent}{0pt}
    \setlength{\parskip}{6pt plus 2pt minus 1pt}}
}{% if KOMA class
  \KOMAoptions{parskip=half}}
\makeatother
\usepackage{xcolor}
\usepackage[margin=1in]{geometry}
\usepackage{color}
\usepackage{fancyvrb}
\newcommand{\VerbBar}{|}
\newcommand{\VERB}{\Verb[commandchars=\\\{\}]}
\DefineVerbatimEnvironment{Highlighting}{Verbatim}{commandchars=\\\{\}}
% Add ',fontsize=\small' for more characters per line
\usepackage{framed}
\definecolor{shadecolor}{RGB}{248,248,248}
\newenvironment{Shaded}{\begin{snugshade}}{\end{snugshade}}
\newcommand{\AlertTok}[1]{\textcolor[rgb]{0.94,0.16,0.16}{#1}}
\newcommand{\AnnotationTok}[1]{\textcolor[rgb]{0.56,0.35,0.01}{\textbf{\textit{#1}}}}
\newcommand{\AttributeTok}[1]{\textcolor[rgb]{0.13,0.29,0.53}{#1}}
\newcommand{\BaseNTok}[1]{\textcolor[rgb]{0.00,0.00,0.81}{#1}}
\newcommand{\BuiltInTok}[1]{#1}
\newcommand{\CharTok}[1]{\textcolor[rgb]{0.31,0.60,0.02}{#1}}
\newcommand{\CommentTok}[1]{\textcolor[rgb]{0.56,0.35,0.01}{\textit{#1}}}
\newcommand{\CommentVarTok}[1]{\textcolor[rgb]{0.56,0.35,0.01}{\textbf{\textit{#1}}}}
\newcommand{\ConstantTok}[1]{\textcolor[rgb]{0.56,0.35,0.01}{#1}}
\newcommand{\ControlFlowTok}[1]{\textcolor[rgb]{0.13,0.29,0.53}{\textbf{#1}}}
\newcommand{\DataTypeTok}[1]{\textcolor[rgb]{0.13,0.29,0.53}{#1}}
\newcommand{\DecValTok}[1]{\textcolor[rgb]{0.00,0.00,0.81}{#1}}
\newcommand{\DocumentationTok}[1]{\textcolor[rgb]{0.56,0.35,0.01}{\textbf{\textit{#1}}}}
\newcommand{\ErrorTok}[1]{\textcolor[rgb]{0.64,0.00,0.00}{\textbf{#1}}}
\newcommand{\ExtensionTok}[1]{#1}
\newcommand{\FloatTok}[1]{\textcolor[rgb]{0.00,0.00,0.81}{#1}}
\newcommand{\FunctionTok}[1]{\textcolor[rgb]{0.13,0.29,0.53}{\textbf{#1}}}
\newcommand{\ImportTok}[1]{#1}
\newcommand{\InformationTok}[1]{\textcolor[rgb]{0.56,0.35,0.01}{\textbf{\textit{#1}}}}
\newcommand{\KeywordTok}[1]{\textcolor[rgb]{0.13,0.29,0.53}{\textbf{#1}}}
\newcommand{\NormalTok}[1]{#1}
\newcommand{\OperatorTok}[1]{\textcolor[rgb]{0.81,0.36,0.00}{\textbf{#1}}}
\newcommand{\OtherTok}[1]{\textcolor[rgb]{0.56,0.35,0.01}{#1}}
\newcommand{\PreprocessorTok}[1]{\textcolor[rgb]{0.56,0.35,0.01}{\textit{#1}}}
\newcommand{\RegionMarkerTok}[1]{#1}
\newcommand{\SpecialCharTok}[1]{\textcolor[rgb]{0.81,0.36,0.00}{\textbf{#1}}}
\newcommand{\SpecialStringTok}[1]{\textcolor[rgb]{0.31,0.60,0.02}{#1}}
\newcommand{\StringTok}[1]{\textcolor[rgb]{0.31,0.60,0.02}{#1}}
\newcommand{\VariableTok}[1]{\textcolor[rgb]{0.00,0.00,0.00}{#1}}
\newcommand{\VerbatimStringTok}[1]{\textcolor[rgb]{0.31,0.60,0.02}{#1}}
\newcommand{\WarningTok}[1]{\textcolor[rgb]{0.56,0.35,0.01}{\textbf{\textit{#1}}}}
\usepackage{longtable,booktabs,array}
\usepackage{calc} % for calculating minipage widths
% Correct order of tables after \paragraph or \subparagraph
\usepackage{etoolbox}
\makeatletter
\patchcmd\longtable{\par}{\if@noskipsec\mbox{}\fi\par}{}{}
\makeatother
% Allow footnotes in longtable head/foot
\IfFileExists{footnotehyper.sty}{\usepackage{footnotehyper}}{\usepackage{footnote}}
\makesavenoteenv{longtable}
\usepackage{graphicx}
\makeatletter
\def\maxwidth{\ifdim\Gin@nat@width>\linewidth\linewidth\else\Gin@nat@width\fi}
\def\maxheight{\ifdim\Gin@nat@height>\textheight\textheight\else\Gin@nat@height\fi}
\makeatother
% Scale images if necessary, so that they will not overflow the page
% margins by default, and it is still possible to overwrite the defaults
% using explicit options in \includegraphics[width, height, ...]{}
\setkeys{Gin}{width=\maxwidth,height=\maxheight,keepaspectratio}
% Set default figure placement to htbp
\makeatletter
\def\fps@figure{htbp}
\makeatother
\setlength{\emergencystretch}{3em} % prevent overfull lines
\providecommand{\tightlist}{%
  \setlength{\itemsep}{0pt}\setlength{\parskip}{0pt}}
\setcounter{secnumdepth}{-\maxdimen} % remove section numbering
\ifLuaTeX
  \usepackage{selnolig}  % disable illegal ligatures
\fi
\usepackage{bookmark}
\IfFileExists{xurl.sty}{\usepackage{xurl}}{} % add URL line breaks if available
\urlstyle{same}
\hypersetup{
  pdftitle={rmd - module2},
  pdfauthor={falsanabriamo},
  hidelinks,
  pdfcreator={LaTeX via pandoc}}

\title{rmd - module2}
\author{falsanabriamo}
\date{2024-06-14}

\begin{document}
\maketitle

\section{This is a level 1 header}\label{this-is-a-level-1-header}

\subsection{R Markdown}\label{r-markdown}

\subsubsection{This is a level 3 header}\label{this-is-a-level-3-header}

Here is a link to \href{https://www.google.com/}{Google}

Here is a word in \textbf{bold} and another word in \textbf{bold}.

Here is a word in \emph{italics} and another word in \emph{italics}.

When we compile our document, we are using the \texttt{rmarkdown}
package.

Here are some example R commands:

\begin{verbatim}
2+2
mean(c(1,2,3,4,5))
\end{verbatim}

\begin{Shaded}
\begin{Highlighting}[]
\DecValTok{2}\SpecialCharTok{+}\DecValTok{2}
\end{Highlighting}
\end{Shaded}

\begin{verbatim}
## [1] 4
\end{verbatim}

\begin{Shaded}
\begin{Highlighting}[]
\FunctionTok{mean}\NormalTok{(}\FunctionTok{c}\NormalTok{(}\DecValTok{1}\NormalTok{,}\DecValTok{2}\NormalTok{,}\DecValTok{3}\NormalTok{,}\DecValTok{4}\NormalTok{,}\DecValTok{5}\NormalTok{))}
\end{Highlighting}
\end{Shaded}

\begin{verbatim}
## [1] 3
\end{verbatim}

\begin{Shaded}
\begin{Highlighting}[]
\CommentTok{\#Cuando no se pone \{r\}, no se ejecuta el comando, solo cambia el formato de texto}
\end{Highlighting}
\end{Shaded}

Here is an example of a non-numbered list:

\begin{itemize}
\tightlist
\item
  Breakfast

  \begin{itemize}
  \tightlist
  \item
    food

    \begin{itemize}
    \tightlist
    \item
      eggs
    \item
      toast
    \item
      bacon
    \end{itemize}
  \item
    drink

    \begin{itemize}
    \tightlist
    \item
      apple juice
    \end{itemize}
  \end{itemize}
\item
  Lunch

  \begin{itemize}
  \tightlist
  \item
    taco
  \end{itemize}
\item
  Dinner

  \begin{itemize}
  \tightlist
  \item
    baked chicken
  \item
    broccoli
  \item
    rice
  \end{itemize}
\end{itemize}

\emph{Recordar que aquí está la regla de los cuatro espacios: cuatro
espacios para una sangría, ocho espacios para dos sangrías\ldots{}}

Here is an example of a numbered list:

\begin{enumerate}
\def\labelenumi{\arabic{enumi}.}
\tightlist
\item
  Breakfast

  \begin{enumerate}
  \def\labelenumii{\alph{enumii}.}
  \tightlist
  \item
    food

    \begin{enumerate}
    \def\labelenumiii{\roman{enumiii}.}
    \tightlist
    \item
      eggs
    \item
      toast
    \item
      bacon
    \end{enumerate}
  \item
    drink

    \begin{enumerate}
    \def\labelenumiii{\roman{enumiii}.}
    \tightlist
    \item
      apple juice
    \end{enumerate}
  \end{enumerate}
\item
  Lunch

  \begin{enumerate}
  \def\labelenumii{\alph{enumii}.}
  \tightlist
  \item
    taco
  \end{enumerate}
\item
  Dinner

  \begin{enumerate}
  \def\labelenumii{\alph{enumii}.}
  \tightlist
  \item
    baked chicken
  \item
    broccoli
  \item
    rice
  \end{enumerate}
\end{enumerate}

Here is an example of a blockquote:

\begin{quote}
This is a block quote. This paragraph has two lines.

\begin{enumerate}
\def\labelenumi{\arabic{enumi}.}
\tightlist
\item
  This is a list inside a block quote.
\item
  Second item.
\end{enumerate}
\end{quote}

Here is an example of a nested blockquote:

\begin{quote}
This is a block quote. This paragraph has two lines.

\begin{quote}
This text is nested
\end{quote}
\end{quote}

Here is an example of code in a blockquote:

\begin{quote}
\begin{verbatim}
2+2
mean(c(1,2,3,4,5))
\end{verbatim}
\end{quote}

\subsection{Gráfica}\label{gruxe1fica}

\includegraphics{rmd_module2_files/figure-latex/pressure-1.pdf}

\begin{Shaded}
\begin{Highlighting}[]
\FunctionTok{plot}\NormalTok{(pressure)}
\end{Highlighting}
\end{Shaded}

\includegraphics{rmd_module2_files/figure-latex/pressure2-1.pdf}

\begin{Shaded}
\begin{Highlighting}[]
\CommentTok{\#tamaño de la figura = 4x4 pulgadas}
\DocumentationTok{\#\#como no se especificó echo=FALSE, se mostrará el código que genera la figura}
\DocumentationTok{\#\#\#estas opciones también pueden cambiarse desde el ícono de engranaje de cada chunk}
\end{Highlighting}
\end{Shaded}

\#\#Tabla

\begin{verbatim}
##   speed dist
## 1     4    2
## 2     4   10
## 3     7    4
## 4     7   22
## 5     8   16
## 6     9   10
\end{verbatim}

\begin{longtable}[]{@{}rr@{}}
\caption{top six rows of cars dataset}\tabularnewline
\toprule\noalign{}
speed & dist \\
\midrule\noalign{}
\endfirsthead
\toprule\noalign{}
speed & dist \\
\midrule\noalign{}
\endhead
\bottomrule\noalign{}
\endlastfoot
4 & 2 \\
4 & 10 \\
7 & 4 \\
7 & 22 \\
8 & 16 \\
9 & 10 \\
\end{longtable}

\#\#Ecuaciones

\[ Y = \beta_0 + \beta1x \] Esta escritura de ecuaciones utiliza el
formato de LaTeX.

\subsection{Insert Images}\label{insert-images}

Here is an image inserted

\begin{figure}
\centering
\includegraphics{sunstar/sunstar.png}
\caption{sunstar}
\end{figure}

Here is the R logo

\begin{figure}
\centering
\includegraphics{https://www.r-project.org/logo/Rlogo.png}
\caption{r logo}
\end{figure}

\subsection{Insert an Animated GIF and
Video}\label{insert-an-animated-gif-and-video}

\begin{figure}
\centering
\includegraphics{sunstar/sunstar.gif}
\caption{sunstar}
\end{figure}

\begin{figure}
\centering
\includegraphics{sunstar/sunstar.mp4}
\caption{sunstar}
\end{figure}

\subsection{Insert text with some
footnotes}\label{insert-text-with-some-footnotes}

Here is a footnote reference \footnote{Here is the footnote.} and
another \footnote{Here's one with multiple blocks.}

Here is an inline footnote \footnote{Inlines notes are easier to write,
  since you don't have to pick an identifier and move down to type the
  note.}

\end{document}
